\documentclass[11pt]{article}

\input{Shared.tex}

\parindent0in
\pagestyle{plain}
\thispagestyle{plain}

\usepackage{csquotes}
\usepackage[shortlabels]{enumitem}

\newcommand{\dated}{\today}
\newcommand{\token}[1]{\langle \text{#1} \rangle}

\begin{document}

\textbf{Introduction to the Theory of
Computation}\hfill\textbf{\myname}\\[0.01in]
\textbf{Chapter 8: Space Complexity}\hfill\textbf{\dated}\\
\smallskip\hrule\bigskip

\begin{problem}{8.13}
Define $A_{\text{LBA}} = \{\langle M,w \rangle \ | \ M \text{ is an LBA that accepts input } w\}$. Show that $A_{\text{LBA}}$ is PSPACE-complete.
\end{problem}

\begin{proof}
To show that $A_{\text{LBA}}$ is PSPACE-complete, we must show that it is in PSPACE and that all PSPACE problems are polynomial time reducible to it. The algorithm given in Theorem 5 decides $A_{\text{LBA}}$\footnote{An LBA operates in exactly $n$ space. This solution applies to both DLBA and NLBA due to Savitch's Theorem.}. This algorithm runs in polynomial space and so $A_{\text{LBA}} \in$ PSPACE. \\

To prove that $A_{\text{LBA}}$ is PSPACE-hard, we give a polynomial time reduction from $GG$ to $A_{\text{LBA}}$. This reduction converts a generalized geography graph $G$ to an LBA $L$ and a string $w$ so that Player I has a winning strategy for the game played on $G$, iff $L$ accepts $w$. The reduction $F$ from $GG$ to $A_{\text{LBA}}$ uses a modified version $L$ of algorithm $M$ given in Theorem 8.14. The algorithm $L$ takes input $\langle G, \ b, \ t \rangle$, where $G$ is a directed graph, $b$ is a node of $G$ and $t$ is a string. The algorithm $L$ operates exactly as algorithm $M$, but only uses space occupied by $t$ as its work memory. Hence, $L$ is an LBA. \\

$F =$ \textquotedblleft On input $\langle G, \ b \rangle$, where $G$ is a directed graph, $b$ is a node of $G$:
\begin{enumerate}
\item Compute the amount of space $n$ required by $M$ on input $\langle G, \ b \rangle$\footnote{Space complexity analysis of algorithm $M$ is given on page 346.}.
\item Construct $L$ as described in proof idea above.
\item Let $w = \langle G, \ b, \ t \rangle$, where $t$ is a string of length $n$.
\item Output $\langle L, \ w \rangle$.\textquotedblright
\end{enumerate}
\end{proof}

\end{document}